% !TeX program = pdflatex
%
% Evidence-Bound Retrieval: A Substrate for CoALA's Episodic Memory Layer
%
% Position paper; companion to the main clinical-AI manuscript
%   "Evidence-Bound Retrieval for Clinical AI" (Anthony, 2026).
%
% Submission targets: NeurIPS 2026 workshops (FMDM, R2-FM),
% ICLR 2027 workshops, ACL position-paper track.
%
% This file uses a generic article class so it can be retargeted to
% specific workshop style files (neurips_2024.sty, iclr2026_conference.sty,
% acl.sty) by swapping the documentclass + \usepackage lines.

\documentclass[11pt,a4paper]{article}

\usepackage[margin=1in]{geometry}
\usepackage[T1]{fontenc}
\usepackage[utf8]{inputenc}
\usepackage{lmodern}
\usepackage{microtype}
\usepackage{amsmath,amssymb,amsthm}
\usepackage{graphicx}
\usepackage{booktabs}
\usepackage{xcolor}
\usepackage{hyperref}
\usepackage{url}
\usepackage{enumitem}
\usepackage{listings}

\hypersetup{
  colorlinks=true,
  linkcolor=blue!50!black,
  citecolor=blue!50!black,
  urlcolor=blue!50!black
}

\lstset{
  basicstyle=\ttfamily\footnotesize,
  breaklines=true,
  showstringspaces=false,
  columns=fullflexible,
  frame=single,
  framesep=4pt,
  framerule=0.4pt,
}

\newtheorem{theorem}{Theorem}
\newtheorem{definition}{Definition}
\newtheorem{remark}{Remark}

\title{Evidence-Bound Retrieval: \\
       A Substrate for CoALA's Episodic Memory Layer}

\author{Ramene Anthony \\
        Independent Researcher \\
        \url{https://github.com/ramene/memory-oracle}}

\date{May 2026}

\begin{document}

\maketitle

\begin{abstract}
The CoALA framework~\cite{sumers2024coala} decomposes a language agent's memory into four
types --- \emph{working}, \emph{semantic}, \emph{procedural}, and \emph{episodic} --- and
identifies episodic memory as ``the hardest of these to get right'' because of the
deletion/obsolescence problem: \emph{what} should the agent forget, and \emph{when}?
Production systems currently address this through \emph{distillation}: decide what is
worth remembering, compress it, and discard the rest. We argue that distillation is the
wrong primitive. The hard question is not selection; it is \emph{precedence}. We present
\textbf{Evidence-Bound Retrieval} (EBR), an alternative episodic-memory substrate that
\emph{never deletes} and \emph{always accretes}: corrections are written as dated
amendment records beside the canonical file and a structural precedence invariant
guarantees they reach retrieval first by construction --- not by similarity, not by a
trained critic, not by reinforcement signal. We characterise EBR formally, show empirical
results from a companion manuscript (N=1{,}000 clinical and trading queries: 100\%
precedence accuracy; a real-corpus probe of 239 documents across 21 projects over 108
days: 6/6 cross-session corrections retrievable), and detail how EBR composes with the
other three CoALA memory types. EBR is offered as a concrete implementation of CoALA's
episodic layer that re-frames CoALA's hardest type from \emph{when to forget} to
\emph{how to amend}.
\end{abstract}

\section{Introduction}
\label{sec:intro}

\noindent
The Princeton CoALA framework~\cite{sumers2024coala} has become a de-facto vocabulary
for describing memory in language-agent systems. In its taxonomy, agent memory
decomposes into four orthogonal types --- \emph{working} (the active context window),
\emph{semantic} (durable knowledge such as project conventions), \emph{procedural} (skill
definitions and step-by-step routines), and \emph{episodic} (the agent's record of past
sessions and what it learned from them).

CoALA is explicit about which of the four is hardest. The framework names
\emph{episodic memory} as the layer for which there is no settled answer:
\begin{quote}
``\emph{Episodic memory is also the hardest type of these to get right because what do
you delete? When does information become obsolete?}''~\cite{ibm2026memory}
\end{quote}
The CoALA paper itself frames episodic memory as ``experience preserved for later use''
and acknowledges that practical systems must decide what to retain and what to discard.
Production systems today operationalise this through \emph{distillation}: each session
emits a compressed summary, the substrate decides what is worth saving (typically via
similarity, recency, or a trained relevance critic), and the rest is dropped.

We argue that this framing makes the problem harder than it needs to be. The
distillation framing assumes the substrate must \emph{choose} between past evidence and
present evidence. The choice is irreversible (you cannot recover what you discarded) and
unaccountable (the choice was made by a heuristic, not by the operator). \textbf{What if
the substrate is not required to choose?}

This paper presents \textbf{Evidence-Bound Retrieval (EBR)}, an episodic-memory substrate
that takes the opposite stance: it \emph{never deletes} and it \emph{always accretes}.
Corrections, updates, and amendments are written as dated, operator-authored records
that sit beside the original canonical document, and a single structural invariant
guarantees they reach retrieval before the canonical body --- not because a model judged
them more relevant, but because the substrate concatenates them first by construction.

EBR re-asks the question. The hardest part of CoALA's episodic layer is not deletion. It
is \emph{precedence}.

\paragraph{Contributions.} This paper makes three contributions.
\begin{enumerate}[leftmargin=2em,itemsep=2pt]
  \item We position EBR explicitly within the CoALA
        framework~\cite{sumers2024coala} as a concrete substrate for its episodic
        memory layer, and we show that EBR is composable (rather than a replacement)
        for the working, semantic, and procedural layers.
  \item We characterise EBR's structural precedence invariant (Section~\ref{sec:ebr}),
        which is the mechanism by which the substrate guarantees that amendments
        win over canonical content at retrieval time --- by file-layout
        construction, not by relevance estimation.
  \item We summarise empirical evidence from a companion manuscript~\cite{anthony2026ebr}
        --- 100\% precedence accuracy on N=1{,}000 clinical and trading queries plus
        a real-corpus probe of the author's own production substrate --- and we
        outline how EBR addresses the open questions Sumers et al.\ identify
        (Section~\ref{sec:composition}).
\end{enumerate}

\section{The Episodic-Memory Problem in CoALA}
\label{sec:problem}

\noindent
Sumers et al.~\cite{sumers2024coala} cast episodic memory as the layer that gives a
language agent \emph{experience}. Working memory is volatile and bounded by the context
window. Semantic memory captures durable facts. Procedural memory captures durable
skills. Episodic memory captures \emph{what happened}: the trajectory of a past session,
the decision the agent made, the outcome it observed, the lesson it should carry forward.

In CoALA's terms a naive episodic implementation stores every session transcript
verbatim. The paper notes this is rarely useful in practice: transcripts are voluminous,
mostly irrelevant on later retrieval, and contain the same information multiple times.
Production systems therefore \emph{distill}: they keep summaries, not transcripts; they
keep ``what was learned,'' not ``what happened.''

Distillation forces three decisions on the substrate, each of which CoALA acknowledges
as unresolved:

\begin{enumerate}[leftmargin=2em,itemsep=2pt]
  \item \textbf{Selection.} Which fragments of a session are worth keeping? The
        substrate must apply a criterion --- typically learned (a trained relevance
        critic, e.g.~\cite{asai2023selfrag,yan2024crag}), or heuristic (recency,
        embedding similarity to a query), or model-judged (an LLM summarises the
        session). Each criterion can drop facts that turn out to matter later.
  \item \textbf{Obsolescence.} When a previously-stored fact stops being true --- the
        user changes jobs, a recommended dosage is corrected, a coding convention is
        retired --- how does the substrate know? In current systems it does not. The
        old fact remains retrievable until something overwrites or removes it, and
        nothing structurally guarantees the new fact wins.
  \item \textbf{Deletion.} If a fact is found to be obsolete, should it be removed?
        If yes, the system has lost evidence of what was once true; if no, both
        versions co-exist with no structural ordering and retrieval picks whichever
        ranks higher under the current scoring function.
\end{enumerate}

These decisions interact. A system that distils aggressively (keeps less) loses
context. A system that distils conservatively (keeps more) reintroduces the transcript
problem and increases the chance of returning superseded content. There is no point on
this trade-off curve at which the substrate is structurally honest about \emph{when} a
fact was true and whether it still is.

\paragraph{What goes missing.} The distillation framing is silent on three properties
that an operator handling consequential decisions (clinical care, algorithmic trading,
production code review) reasonably expects from a memory layer:

\begin{itemize}[leftmargin=2em,itemsep=2pt]
  \item \emph{Auditability.} Which prior version was current at the time the agent
        acted? A distilled summary does not preserve the timeline.
  \item \emph{Authorship.} Who corrected the prior assertion? A distilled summary
        does not preserve the operator's signature on the correction.
  \item \emph{Precedence.} If two contradictory facts are present, which wins at
        retrieval, and is the answer the same on every retrieval? Under similarity-
        or critic-based retrieval, the answer can vary with the query phrasing.
\end{itemize}

\noindent
We argue that the right substrate does not choose between facts. It \emph{stacks} them
in operator-authored, dated order, and it makes the most recent stack-element retrievable
first by construction. That is the EBR proposal.

\section{Evidence-Bound Retrieval}
\label{sec:ebr}

\noindent
Evidence-Bound Retrieval (EBR) is a property of a retrieval substrate, not an algorithm
on top of one. The property is:

\begin{definition}[Evidence binding]
\label{def:ebr}
A retrieval system is \emph{evidence-bound} if, for every retrieved fragment $f$ at time
$t$, $f$ is the most recent operator-authored evidence regarding its subject as of $t$,
by construction --- i.e.\ the property is enforced by the substrate's file layout and
read order, not by a similarity score, a learned relevance critic, or a reinforcement
signal.
\end{definition}

\noindent
We describe one realisation of an EBR substrate. The realisation is small (it is a
file-layout convention plus a deterministic merge), and the smallness is part of the
argument: an episodic-memory substrate need not be a model.

\subsection{Amendment records}
\label{sec:amendments}

\noindent
A canonical document is a plain Markdown file with optional YAML front-matter. An
\emph{amendment record} is a JSONL sidecar file sitting beside the canonical file:

\begin{lstlisting}[language=,caption={Amendment record layout.},captionpos=b]
docs/examples/clinical-records/patient-jane/memory/
  medication_anticoagulant.md
  medication_anticoagulant.md.amendments.jsonl
\end{lstlisting}

Each line of the sidecar is a JSON object recording one correction event:

\begin{lstlisting}[language=,caption={Single amendment event (illustrative; full schema in companion paper).},captionpos=b]
{
  "amended_at":         "2026-03-14T11:02:00Z",
  "amended_by":         "Dr. Reyes, MD",
  "superseded_assertion": "Patient is on warfarin 5 mg/day.",
  "corrected_assertion":  "Patient transitioned to apixaban 5 mg twice daily on 2026-03-12.",
  "live_evidence":      "EHR/encounter/E-71412/note-2.txt#L42",
  "operator_confirmed": true,
  "retention_policy":   "indefinite"
}
\end{lstlisting}

\noindent
Three properties of this layout matter for the rest of the paper:

\begin{enumerate}[leftmargin=2em,itemsep=2pt]
  \item \textbf{Canonical-free.} The canonical Markdown file is never edited. The
        original assertion (``warfarin 5 mg/day'') remains exactly as authored.
        This preserves the audit trail.
  \item \textbf{Append-only.} The sidecar is JSONL: one event per line, written
        last-wins. Older amendments are never overwritten; a re-amendment is a new
        line.
  \item \textbf{Locality.} The sidecar is in the same directory as the canonical
        file with a deterministic extension. No external index is required to find
        the corrections that apply to a given file.
\end{enumerate}

\subsection{The precedence invariant}
\label{sec:precedence}

\noindent
The mechanism by which an amendment beats the canonical body at retrieval is a single
invariant on how the substrate \emph{reads} the file pair. We state it as a theorem.

\begin{theorem}[Structural precedence]
\label{thm:precedence}
Let $C$ be a canonical document and $A$ its associated amendment sidecar with $|A| \geq 0$
records. For every retrieval that returns content from the pair $(C, A)$, the substrate
constructs the returned text $R$ as:
\[
R \;=\; \mathit{block}(A_{\downarrow t}) \;\Vert\; C
\]
where $A_{\downarrow t}$ denotes the amendment events in $A$ ordered by descending
$\mathit{amended\_at}$, $\mathit{block}(\cdot)$ formats them into a leading
``AMENDMENTS'' section, and $\Vert$ is text concatenation.
The amendments precede the canonical body by construction. The order is independent of
similarity score, relevance score, query phrasing, model temperature, and reinforcement
signal.
\end{theorem}

\begin{proof}[Proof sketch]
The substrate's merge routine ($\mathit{merge\_canonical\_with\_amendments}$ in the
reference implementation~\cite{anthony2026ebr}) reads $A$ into memory, sorts events
descending by timestamp, formats them as the leading block, and concatenates $C$. No
ranking step is interposed between sort and concatenation. Therefore for any
$(C, A, t)$ the substrate emits the same $R$, and the amendments precede $C$ in $R$.
\qed
\end{proof}

\noindent
The invariant is intentionally trivial. It is the basis on which EBR claims to be a
\emph{substrate property} rather than a \emph{model behaviour}: the property is read off
the merge code in a single inspection, not validated by ablation studies.

\subsection{Three-tier retrieval vocabulary}
\label{sec:tiers}

\noindent
A practical EBR substrate exposes three distinct retrieval surfaces, all of which respect
the precedence invariant. The tiers differ in latency and in the kind of query they
answer:

\begin{description}[leftmargin=2em,itemsep=4pt]
  \item[Tier 1 (BM25 keyword search).] Sub-second full-text retrieval over the
        canonical corpus, with amendments stitched in at merge time. Suitable for
        ``find me the current value of X.''
  \item[Tier 1.5 (structural surface map).] SQL queries against the substrate's
        index that exploit file-system structure (project, file, line, amendment-
        count). Suitable for ``which files in project P have been amended in the
        last 30 days?''
  \item[Tier 2 (forensic JSONL scan).] Direct scan of amendment sidecars to recover
        the full timeline of corrections on a specific file. Slower but
        exhaustive --- needed for audit, never needed for inference-time retrieval.
\end{description}

\noindent
The three tiers cover the typical question shapes a CoALA-style agent asks of its
episodic layer (current-state lookup, surface-level enumeration, full-history audit)
without forcing any one tier to do all three jobs.

\section{Empirical Evidence}
\label{sec:empirical}

\noindent
This position paper does not introduce new experiments; the experiments live in the
companion manuscript~\cite{anthony2026ebr}. We summarise the headline results so the
positioning argument can be read in isolation.

\paragraph{Synthetic vault stress test (N=1{,}000 queries).}
Two synthetic patient/trader vaults were constructed: a clinical record with a
warfarin~$\to$~apixaban medication amendment and a trading-strategy record with a hard
no-shorting rule added after a real $\$206$ loss event. Each was probed with 500
phrasings of the same retrieval intent (``what is the current X?'' / ``how does the
agent treat Y?''). On the clinical vault, EBR returned the amended assertion as
precedence-correct on \textbf{100.0\%} of queries; vector-RAG returned the canonical
(pre-amendment) assertion on \textbf{90\%} (the canonical document is the higher-
similarity match because it is longer); a control LLM with no retrieval scored
\textbf{0.0\%}. A required-litmus measure --- the fraction of queries on which the
substrate returns the post-amendment assertion --- reads $\mathrm{LLM} = 0.0$,
$\mathrm{Vector} = 0.1$, $\mathrm{EBR} = 1.0$, gap $= 0.9$. The trading-vault numbers
are equivalent.

\paragraph{Real-corpus probe.}
The author's own production substrate --- 239 documents across 21 projects, 108 days
of writing, an unmodified ``live'' EBR system used for normal work --- was probed for
six known cross-session corrections (e.g.\ rename events, error-message corrections,
operator-preference changes). All 6/6 corrections were retrievable in Tier-1 BM25 search
with median latency $257$~ms. This is offered as evidence that the substrate behaves
on operator-authored content the way the synthetic vaults predict --- i.e.\ that the
100\% precedence-accuracy number is not an artifact of synthetic-vault design.

\paragraph{Index hygiene and contention.}
Concurrent write tests (30 amendment events emitted simultaneously across 5 substrates)
returned 30/30 amendments indexed and retrievable, with 7/30 events incurring transient
SQLite-busy errors that the substrate retried internally. No amendment was lost.

\paragraph{The methodological objection.}
A 100\% precedence-accuracy number invites the reviewer's reasonable concern that the
experiment is rigged. The companion manuscript devotes a methodological note to this:
EBR's correctness is \emph{structural}, not \emph{statistical}. The substrate's merge
routine prepends amendments by construction (Theorem~\ref{thm:precedence}). The 100\%
number is therefore expected --- it would be the surprise if it were anything else. The
useful measurement is not whether EBR is 100\% accurate (it is, by construction); it is
whether the comparison baselines (LLM, vector-RAG) are below 100\%, by how much, and
whether the gap is operationally meaningful. They are 0\% and 10\%, respectively, and
the gap is $0.9$ on the litmus.

\section{Composition with CoALA's Other Memory Types}
\label{sec:composition}

\noindent
EBR is offered as an episodic-layer substrate, not a replacement for the rest of CoALA.
This section describes how EBR composes with the working, semantic, and procedural
layers.

\subsection{Semantic $\oplus$ EBR}
\label{sec:semantic-plus-ebr}

\noindent
Sumers et al.\ describe semantic memory as the durable knowledge layer --- in production
systems this is realised as Markdown documents (\texttt{CLAUDE.md}, project READMEs,
operator-curated knowledge bases). EBR does not replace the semantic layer; it
\emph{amends} it. An amendment record can attach to any semantic-layer Markdown file:
the canonical document remains the authoritative knowledge base, and the amendment is
the operator's correction. The pair $(C, A)$ is read into context together, with $A$
preceding by Theorem~\ref{thm:precedence}.

This composition resolves CoALA's deletion problem on the semantic layer without
modifying it: the operator never deletes a fact from \texttt{CLAUDE.md}; the operator
writes an amendment that supersedes the relevant assertion. The canonical document
remains diff-able, blame-able, and reviewable.

\subsection{Procedural $\oplus$ EBR}
\label{sec:procedural-plus-ebr}

\noindent
Procedural memory in CoALA is realised in production as the agent-skills standard
(\texttt{skill.md} files in skill folders~\cite{anthropic2024skills}). A skill is itself
a Markdown document: it can be amended in exactly the same way as a semantic-layer
document. The interesting case is a \emph{procedural correction}: ``the skill used to
say X; under the new tool, do Y instead.'' EBR captures this as an amendment record on
the skill's \texttt{skill.md} without modifying the canonical instructions.

This is the closest EBR comes to giving a language agent something that resembles
classical reinforcement learning's policy update: the skill's effective behaviour
changes (the agent does Y instead of X), but the change is operator-authored, dated,
and reversible by deleting one JSONL line --- not learned, not opaque, not
permanently baked into a policy network.

\subsection{Working $\oplus$ EBR}
\label{sec:working-plus-ebr}

\noindent
Working memory in CoALA is the context window. EBR delivers retrieved content into the
context window the way any RAG-style substrate does: a tool call returns text, the
agent's runtime concatenates it into the prompt. The difference is the text itself ---
because amendments are prepended at merge time, the working memory the agent sees
already has the most recent operator-authored correction at the top.

A practical engineering implication: amendment-merged retrievals are longer than the
canonical-only document (typical observation in the companion paper:
$\approx\!10$--$30$\,KB merged vs.\ $\approx\!2$--$5$\,KB canonical-only). The context-
window budget therefore matters more under EBR than under bare RAG. The companion paper
pairs EBR with FP8 KV-cache compression to keep per-token inference cost
parity~\cite{turboquant2026}; that engineering pairing is orthogonal to the EBR property
itself.

\subsection{What EBR does not replace}

\noindent
EBR does not replace distillation; it complements it. Distillation remains useful for
\emph{compressing} an experience into a remembered shape (``last time we debugged the
auth module, the issue was in the middleware layer'' is a distilled fact, not a
transcript). EBR's contribution is to make distilled facts \emph{amendable}: when a
later session discovers the auth issue was actually in the OAuth callback, the operator
writes an amendment on the distilled note, and from that retrieval onward the agent
sees the corrected version first.

In CoALA terms: distillation answers \emph{what is worth remembering}; EBR answers
\emph{what should still be believed}.

\section{Limitations and Open Problems}
\label{sec:limits}

\paragraph{Multi-author amendment provenance.}
The reference implementation assumes a single operator authors amendments. In multi-
user settings (multiple clinicians on one patient record, multiple traders on one
strategy file), the amendment record needs richer provenance: who authored, who
co-signed, who reviewed. The JSON schema supports the additional fields trivially but
the substrate's retrieval-time merge currently treats all amendments equally; a
\emph{multi-author precedence policy} is needed.

\paragraph{Cryptographic signing.}
Amendments are append-only on the filesystem but not cryptographically signed. A
malicious or accidental edit to an amendment line can change what the substrate
retrieves without leaving evidence beyond the file-system mtime. Pairing the JSONL
sidecar with a Merkle log~\cite{shamir1979} or git-crypt~\cite{gitcrypt2021} signing
chain is a natural next step; the file-layout convention is unaffected.

\paragraph{Cross-substrate amendment portability.}
The amendment sidecar is a file-layout convention, not a wire protocol. Two distinct
substrates that both adopt the convention can read each other's amendments by file
exchange, but there is no agreed-upon schema version field, no migration path between
schema revisions, and no documented mapping to existing standards (FHIR addenda for
clinical~\cite{fhir2023,addendum2018}, trade-confirmation amendments for finance). A
JSON-schema specification with stable version numbering is a near-term need.

\paragraph{Index hygiene for amendment chains.}
A canonical file with $k$ amendments produces a retrieval of $O(k)$ tokens longer than
its bare form. For files amended hundreds of times, the merged retrieval can crowd the
working-memory budget. A garbage-collection policy --- e.g.\ ``compress amendments older
than $T$ into a single rollup record, retaining the JSONL log on disk for audit but not
loading it into context'' --- is sketched in the companion paper but not yet measured
under load.

\paragraph{What this paper does not claim.}
EBR is not claimed to be the only viable episodic-memory substrate, nor to outperform
learned distillation on tasks that genuinely need compression. It is claimed to be a
substrate that \emph{structurally guarantees} the property in Definition~\ref{def:ebr},
which the distillation framing does not.

\section{Related Work}
\label{sec:related}

\noindent
\textbf{Retrieval-time correction (RAG family).} CRAG~\cite{yan2024crag} introduces a
trained retrieval critic that decides whether retrieved documents are reliable;
Self-RAG~\cite{asai2023selfrag} learns to retrieve, generate, and critique inline;
FLARE~\cite{jiang2023flare} actively re-retrieves during generation when the model's
confidence drops. All three operate within the distillation framing --- the substrate
returns what its retriever produces, and a learned component decides what to do with
it. EBR's precedence is structural rather than learned, which makes it cheaper at
inference time (no critic forward pass) and auditable (the precedence decision is
visible in the merge code).

\noindent
\textbf{Long-term memory systems.} MemGPT~\cite{memgpt2023} treats the LLM as an
operating system and pages between a small in-context working set and a larger external
store; LangMem~\cite{langmem2024} provides long-term memory for LangChain agents. Both
address \emph{capacity} (more memory than the context window holds). Neither addresses
\emph{precedence}: when the paged-in memory contradicts a more recent fact, the system
has no structural mechanism to ensure the more recent fact wins.

\noindent
\textbf{Clinical and financial amendment conventions.} The EHR-addendum
literature~\cite{addendum2018} and the FHIR R5 standard~\cite{fhir2023} both establish
that medical records should be amended rather than overwritten, for the same audit-trail
reasons EBR claims. EBR can be read as porting this clinical-record convention to the
general agent memory substrate.

\noindent
\textbf{The CoALA framework itself.} Sumers et al.~\cite{sumers2024coala} are the
proximate cause of this paper: they named the four memory types, articulated the
episodic-layer problem, and provided the vocabulary in which we can ask ``what would a
production-grade episodic-memory substrate look like?'' This paper is a direct response
to that question.

\section{Conclusion}
\label{sec:conclusion}

\noindent
CoALA's four-memory taxonomy will likely outlast any specific implementation built
inside it. Within that taxonomy, episodic memory is identified as the hardest layer
because the deletion/obsolescence problem has no clean answer under a distillation
framing. We have argued that the right move is not a better answer to that question
but a different question. EBR replaces \emph{when to forget} with \emph{how to amend}:
the substrate never deletes, always accretes, and binds retrieval to the most recent
operator-authored evidence by a single structural invariant that is read off the merge
code in one inspection.

The reference implementation~\cite{anthony2026ebr} is small (a file-layout convention
plus a deterministic merge), and the smallness is the point. An episodic-memory
substrate that works under CoALA need not be a model. It can be a discipline about how
files are laid out and read.

\paragraph{What we would like to see next.}
We would welcome (a) an empirical comparison of EBR against learned-distillation
baselines on a public benchmark that includes operator-authored corrections, (b)
adoption of the amendment-record convention by other CoALA-aligned substrates, and (c)
extensions of the precedence invariant to handle multi-author provenance and
cryptographic signing.

\bigskip
\noindent
\textbf{Reproducibility.} The reference implementation, synthetic vaults, and exact
queries used in the empirical results section live at
\url{https://github.com/ramene/memory-oracle}. The companion full manuscript and its
empirical-evaluation notebook are in the same repository under \texttt{paper/lncs/} and
\texttt{notebooks/memory-oracle/}.

\bibliographystyle{plain}
\bibliography{references}

\end{document}
